\documentclass[11pt,a4paper]{article}

% ---------------------------------------------------------
% PACKAGES
% ---------------------------------------------------------
\usepackage[utf8]{inputenc}
\usepackage[T1]{fontenc}
\usepackage{lmodern}
\usepackage{amsmath, amssymb, amsfonts, bm}
\usepackage{geometry}
\usepackage{graphicx}
\usepackage{hyperref}
\usepackage{authblk}
\usepackage{cite}
\usepackage{physics}
\usepackage{mathtools}
\usepackage{tensor}
\usepackage{enumitem}

\geometry{margin=2.4cm}

% ---------------------------------------------------------
% TITLE
% ---------------------------------------------------------
\title{\textbf{Companion VI (v49.1)}\\[0.2cm]
\textbf{Relaxation Geometry and Infrared Structure}\\[0.2cm]
\large Foundations of Relaxation-Driven Cyclic Cosmology (RDCC)}

\author[1]{Michael Lehmann}
\affil[1]{\small Independent Researcher, Relaxation-Driven Cyclic Cosmology (RDCC)\\
\texttt{mi.lehmann@gmx.de}}

\date{September 2026}

\begin{document}
\maketitle

% ========================================================
% ABSTRACT
% ========================================================

\begin{abstract}
This Companion develops the geometric interpretation of relaxation in RDCC.
The relaxation parameter $\alpha(t)$ is a global infrared quantity arising from
the ultraviolet-suppressed dimension-8 operator $T^{(+)}_{\mu\nu}T^{(-)\mu\nu}$.
Relaxation is not a sectoral process but a geometric response of the global
metric. We derive the geometric origin of $\alpha(t)$, establish its infrared
fixpoint $\alpha_{\mathrm{IR}}$, and explain how relaxation modulates scalar
and tensor perturbations. This Companion integrates the ultraviolet sector
(Companion X), the relaxation sector (Companion III), and the tensor-relaxation
coupling (Companion XI).
\end{abstract}

\tableofcontents
\newpage
% ========================================================
% 1 INTRODUCTION
% ========================================================

\section{Introduction}

Relaxation in RDCC is not a thermodynamic process. It is a geometric response
of the global metric to the ultraviolet-suppressed coupling between the visible
and mirror sectors. The relaxation parameter $\alpha(t)$ is globally defined
and arises from the inter-sector operator
\begin{equation}
\mathcal{O}_{\mathrm{int}} = T^{(+)}_{\mu\nu} T^{(-)\mu\nu}.
\end{equation}

In earlier versions (up to v48.0), relaxation was interpreted as a sectoral
process. In v49.1, this interpretation is corrected: relaxation is a global
geometric effect that modulates scalar and tensor perturbations and saturates
in the infrared due to RG flow.

This Companion develops the geometric structure of relaxation in a form that is
mathematically precise, physically transparent, and fully compatible with the
global structure of RDCC.
% ========================================================
% 2 GEOMETRIC ORIGIN OF RELAXATION
% ========================================================

\section[Geometric Origin of alpha(t)]{Geometric Origin of $\alpha(t)$}

Relaxation arises from the geometric contraction of the global metric with the
dimension-8 operator,
\begin{equation}
\mathcal{O}_{\mathrm{int}} = T^{(+)}_{\mu\nu} T^{(-)\mu\nu}.
\end{equation}

The effective action contains
\begin{equation}
S_{\mathrm{int}}
=
\int d^4x \sqrt{-g} \,
\frac{1}{M^4} \,
T^{(+)}_{\mu\nu} T^{(-)\mu\nu}.
\label{eq:dim8action}
\end{equation}

The relaxation parameter is defined as the geometric response of the global
metric to this operator,
\begin{equation}
\alpha(t) = \frac{1}{M^4} \, \mathcal{G}(t),
\end{equation}
where $\mathcal{G}(t)$ encodes the infrared geometry of the global state.

\subsection[Global Definition]{Global Definition}

Relaxation is globally defined:
\begin{equation}
\alpha(t) = \alpha_{\mathrm{global}}(t).
\end{equation}

The visible and mirror sectors measure $\alpha(t)$ only through projection.
% ========================================================
% 3 RG FLOW AND INFRARED FIXPOINT
% ========================================================

\section[RG Flow and Infrared Fixpoint]{RG Flow and Infrared Fixpoint}

The geometric response saturates in the infrared due to RG flow,
\begin{equation}
\dot{\alpha}(t) = - \Gamma_{\mathrm{IR}} \, \alpha(t),
\end{equation}
with solution
\begin{equation}
\alpha(t) = \alpha_{\mathrm{IR}}
+
\left( \alpha_0 - \alpha_{\mathrm{IR}} \right)
e^{-\Gamma_{\mathrm{IR}} t}.
\end{equation}

The infrared fixpoint is
\begin{equation}
\alpha_{\mathrm{IR}} = \lim_{t \to \infty} \alpha(t).
\end{equation}

This fixpoint controls all deviations from $\Lambda$CDM:
\begin{align}
n_s - 1 &\sim -2\alpha_{\mathrm{IR}}, \\
\Delta N_{\mathrm{eff}} &\sim \alpha_{\mathrm{IR}}, \\
\Delta(f\sigma_8) &\sim \alpha_{\mathrm{IR}}, \\
\Omega_{\mathrm{GW}} &\propto \alpha_{\mathrm{IR}}^2.
\end{align}
% ========================================================
% 4 GEOMETRIC MODULATION OF PERTURBATIONS
% ========================================================

\section{Geometric Modulation of Perturbations}

Relaxation modulates both scalar and tensor perturbations. The global metric
responds to the inter-sector coupling, and this response is encoded in
$\alpha(t)$.

\subsection[Scalar Modes]{Scalar Modes}

Scalar perturbations receive a correction proportional to $\alpha(t)$,
\begin{equation}
\Phi(k,t) = \Phi_{\Lambda\mathrm{CDM}}(k,t)
\left[ 1 + \alpha(t) \right].
\end{equation}

\subsection[Tensor Modes]{Tensor Modes}

Tensor modes receive a logarithmic modulation and quadratic suppression,
\begin{equation}
\Omega_{\mathrm{GW}}(k)
=
\Omega_{\mathrm{GW}}^{(0)}(k)
\left[ 1 + \alpha_{\mathrm{IR}} \log(k/k_0) \right],
\end{equation}
with
\begin{equation}
\Omega_{\mathrm{GW}} \propto \alpha_{\mathrm{IR}}^2.
\end{equation}

These signatures are unique to RDCC and arise from the global coherence of
tensor modes.
% ========================================================
% 5 CONCLUSION
% ========================================================

\section{Conclusion}

This Companion develops the geometric interpretation of relaxation in RDCC.
Relaxation is a global infrared effect arising from the ultraviolet-suppressed
dimension-8 operator $T^{(+)}_{\mu\nu}T^{(-)\mu\nu}$. The relaxation parameter
$\alpha(t)$ is globally defined, saturates in the infrared due to RG flow, and
controls all deviations from $\Lambda$CDM.

These results integrate the relaxation sector (Companion III), the ultraviolet
sector (Companion X), the tensor-relaxation coupling (Companion XI), and the
Boltzmann hierarchy (Companion VII). Together, they establish the geometric
foundation of RDCC’s infrared physics.

\end{document}
